\documentclass[aspectratio=169,10pt]{beamer}
\usetheme{quiet}

\usepackage{amsmath,amssymb}
\usepackage{booktabs}
\usepackage{colortbl}
\usepackage{eurosym}
\usepackage{pgfplots}
\pgfplotsset{compat=1.16}
\usepgfplotslibrary{groupplots}

\pgfplotsset{quietaxis/.style={
  axis lines=left, axis line style={hair, line width=0.6pt},
  tick style={draw=none}, tick label style={font=\scriptsize, color=quiet},
  label style={font=\scriptsize, color=quiet},
  title style={font=\small, color=ink, yshift=-0.4em},
  every axis plot/.append style={line width=1.4pt},
  legend style={draw=none, fill=none, font=\scriptsize, text=ink},
  legend cell align=left}}

\title{Use It or Lose It?}
\subtitle{Programmable money, negative rates and the design of the digital euro}
\author{[Presenter name]}
\date{October 2026}

\begin{document}

% ==================================================================
\begin{frame}[plain]
  \titlepage
\end{frame}

% 2 ================================================================
\begin{frame}{Monetary policy stalls when money is hoarded and rates cannot fall below zero}{The problem}
  \vspace{0.5cm}
  \begin{columns}[T]
    \begin{column}{0.44\textwidth}
      \colhead{Money is not spent}
      \bignum{accent}{\euro 8.8tn}{Eurosystem balance sheet at its peak, June 2022}
      \vspace{0.5em}
      \quiet{\small yet inflation averaged 0.3\% in 2014--16: new money was held, not spent}
    \end{column}
    \begin{column}{0.44\textwidth}
      \colhead{Rates cannot fall}
      \bignum{accent}{8 years}{with the deposit facility rate at or below zero, 2014--22}
      \vspace{0.5em}
      \quiet{\small floor at $-0.50\%$: below that, savers switch to cash, which pays zero}
    \end{column}
  \end{columns}
  \vspace{0.8cm}
  \begin{center}
    $M\,\hl{V}\!\downarrow \;=\; P\,Y$ \qquad\qquad $i \;\ge\; -\,c_{\text{cash}}$
  \end{center}
\end{frame}

% 3 ================================================================
\begin{frame}{Vouchers with a deadline get money spent, and a programmable digital euro could become one}{A tempting fix}
  \vspace{0.5cm}
  \begin{columns}[T]
    \begin{column}{0.44\textwidth}
      \colhead{China, digital coupons, 2020}
      \bignum{accent}{\textyen 3}{extra spending per \textyen 1 of subsidy}
    \end{column}
    \begin{column}{0.44\textwidth}
      \colhead{France, randomised prepaid cards}
      \bignum{accent}{61\% \quiet{\normalsize vs} \quiet{23\%}}{one-month MPC with vs without a three-week expiry}
    \end{column}
  \end{columns}
  \vspace{0.9cm}
  \takeaway{A digital euro is programmable. Could money itself carry a deadline?}
  \source{Xing, Zou, Yin, Wang \& Li (AEJ: Macro 2023); Ding, Jiang, Msall \& Notowidigdo (2024); Boehm, Fize \& Jaravel (2024).}
\end{frame}

% 4 ================================================================
\begin{frame}{But an expiry date pushes out more spending, ends in a fire sale and splits money into vintages}{The flaw}
  \begin{columns}[T]
    \begin{column}{0.30\textwidth}
      \vspace{0.6em}
      \colhead{Model}
      \small
      A euro is worth $V$ to its holder. A good purchase is worth $1$, a fire sale $\phi<1$.
      Searching at speed $\kappa$ costs $\kappa^2/2a$, so
      \[
        \kappa^* = a\,(1-V)
      \]
    \end{column}
    \begin{column}{0.66\textwidth}
      \begin{tikzpicture}
        \begin{groupplot}[group style={group size=2 by 1, horizontal sep=1.2cm},
            quietaxis, width=0.56\textwidth, height=3.7cm, xmin=0, xmax=31, xlabel={day}, yticklabel style={/pgf/number format/fixed, /pgf/number format/precision=2}]
          \nextgroupplot[title={Velocity}, ymin=0.03, ymax=0.11, ytick={0.04,0.06,0.08,0.10},
              legend style={at={(0.02,0.98)}, anchor=north west}]
            \addplot[warm] table[x=day, y=kappa_expiry, col sep=comma] {../figures/expiry_paths.csv};
            \addlegendentry{expiry date}
            \addplot[accent] table[x=day, y=kappa_decay, col sep=comma] {../figures/expiry_paths.csv};
            \addlegendentry{continuous decay}
            \addplot[only marks, mark=*, mark size=1.6pt, warm] coordinates {(30,0.10)};
          \nextgroupplot[title={Value of one unit}, ymin=0.45, ymax=0.85, ytick={0.5,0.6,0.7,0.8}]
            \addplot[warm] table[x=day, y=value_expiry, col sep=comma] {../figures/expiry_paths.csv};
            \addplot[accent] table[x=day, y=value_decay, col sep=comma] {../figures/expiry_paths.csv};
            \addplot[quiet, line width=0.6pt, dashed] coordinates {(0,0.5) (31,0.5)};
            \node[font=\scriptsize, color=quiet, anchor=south west] at (axis cs:0.5,0.5) {fire-sale value $\phi$};
        \end{groupplot}
      \end{tikzpicture}
    \end{column}
  \end{columns}
  \vspace{0.25cm}
  \begin{columns}[T]
    \begin{column}{0.30\textwidth}\centering
      {\Large\color{warm}84\%} {\small vs} {\Large\color{accent}73\%}\\[-0.1em]{\scriptsize\quiet{useful spending within 30 days}}
    \end{column}
    \begin{column}{0.30\textwidth}\centering
      {\Large\color{warm}16\%}\\[-0.1em]{\scriptsize\quiet{of issuance dumped on the last day}}
    \end{column}
    \begin{column}{0.30\textwidth}\centering
      {\Large\color{warm}0.80 $\rightarrow$ 0.50}\\[-0.1em]{\scriptsize\quiet{value of one euro, by its age}}
    \end{column}
  \end{columns}
  \source{Both schemes calibrated to the same average velocity over 30 days; $a=0.2$, $\phi=0.5$, $r=4\%$.}
\end{frame}

% 5 ================================================================
\begin{frame}{Continuous decay keeps every euro equal: it is simply a negative rate that reaches every wallet}{The repair}
  \vspace{0.4cm}
  \begin{columns}[T]
    \begin{column}{0.48\textwidth}
      \small
      \begin{tabular}{@{}lll@{}}
        & \textcolor{warm}{Expiry} & \textcolor{accent}{Decay} \\
        \arrayrulecolor{hair}\midrule
        Fire sale at the deadline & 3--55\% & none \\
        Value of one unit & by age & equal \\
        One euro $=$ one euro & \textcolor{warm}{no} & yes \\
        Applies to all money & no & yes \\
        Is it an interest rate? & no & $i=-\lambda$ \\
      \end{tabular}
    \end{column}
    \begin{column}{0.44\textwidth}
      \colhead{Velocity under decay}
      \[
        \kappa^* \approx \sqrt{2a\,(r+\lambda)}
      \]
      \quiet{\small the square-root rule of Baumol--Tobin money demand}
      \vspace{0.8em}

      \small
      \begin{tabular}{@{}lccc@{}}
        Decay $\lambda$ & 1\% & 2\% & 3\% \\
        \arrayrulecolor{hair}\midrule
        Velocity & $+12\%$ & $+22\%$ & $+32\%$ \\
      \end{tabular}
    \end{column}
  \end{columns}
  \vspace{0.7cm}
  \takeaway{It unblocks both channels: money moves faster, and the rate reaches retail balances.}
\end{frame}

% 6 ================================================================
\begin{frame}{Yet the digital euro is set at zero, which makes it cash without storage costs and raises the lower bound}{The twist}
  \vspace{0.15cm}
  \begin{center}
    {\large\itshape ``The digital euro shall not bear interest.''}\\[0.2em]
    {\scriptsize\quiet{Draft regulation on the digital euro, Art.\ 16(8)}}
  \end{center}
  \vspace{0.2cm}
  \begin{columns}[c]
    \begin{column}{0.52\textwidth}
      \begin{tikzpicture}[x=5.2cm]
        \draw[hair, line width=0.8pt] (-0.75,0) -- (0.45,0);
        \foreach \x/\l in {-0.5/{$-0.5\%$}, 0/{$0$}} \draw[quiet] (\x,0.08) -- (\x,-0.08) node[below, font=\scriptsize, color=quiet] {\l};
        \fill[warm] (-0.5,0) circle (2.6pt);
        \node[above, font=\scriptsize, color=warm, align=center] at (-0.5,0.15) {cash\\floor $-c$};
        \fill[accent] (0,0) circle (2.6pt);
        \node[above, font=\scriptsize, color=accent, align=center] at (0,0.15) {zero-rate\\digital euro};
        \draw[->, accent, line width=1pt] (-0.46,-0.55) -- (-0.04,-0.55);
        \node[below, font=\scriptsize, color=accent] at (-0.25,-0.6) {the lower bound moves up};
        \node[font=\scriptsize, color=quiet] at (0.32,-0.25) {deposit rate};
      \end{tikzpicture}
    \end{column}
    \begin{column}{0.42\textwidth}
      \small
      Banks cannot pay less than the best central-bank alternative:
      \[
        r_D \;\ge\; \max_j\,\{\, i_j - c_j \,\}
      \]
      Cash: $0-c<0$. \quad Digital euro: $0-0=0$.
    \end{column}
  \end{columns}
  \vspace{0.3cm}
  \takeaway{The lower bound and disintermediation are one problem: a central-bank rival to deposits.}
\end{frame}

% 7 ================================================================
\begin{frame}{A hard holding limit stops runs in a crisis, but blocks ordinary users in normal times}{The trade-off}
  \vspace{0.5cm}
  \begin{columns}[T]
    \begin{column}{0.44\textwidth}
      \colhead{Normal times: money for payments}
      \bignum{warm}{18\%}{of people want to hold more than \euro 3{,}000}
      \vspace{0.6em}
      \quiet{\small A cap blocks them. A price would only make them pay.}\\[0.5em]
      \hl{\small $\Rightarrow$ price is enough}
    \end{column}
    \begin{column}{0.44\textwidth}
      \colhead{Crisis: flight to safety}
      \bignum{accent}{\euro 4.7tn}{deposit outflow without a cap \quiet{(\euro 0.65tn with \euro 3{,}000)}}
      \vspace{0.6em}
      \quiet{\small A rate gap of 2\% means nothing next to the fear of losing a deposit.}\\[0.5em]
      \hl{\small $\Rightarrow$ only quantity works}
    \end{column}
  \end{columns}
  \source{Survey demand: Lambert et al.\ (ECB WP 2980, 2024). Outflows: own calibration to ECB technical data on holding limits (Oct 2025).}
\end{frame}

% 8 ================================================================
\begin{frame}{When costs have a cliff and benefits do not, price and quantity together beat either alone}{The principle}
  \begin{columns}[c]
    \begin{column}{0.52\textwidth}
      \begin{tikzpicture}
        \begin{axis}[quietaxis, width=\textwidth, height=5.2cm, xmin=0, xmax=10, ymin=0, ymax=6,
            xtick={2,7}, xticklabels={$X_0$, $L$}, ytick={1.6}, yticklabels={$e$},
            xlabel={outflow from deposits $X$}, ylabel={marginal value / cost}]
          \addplot[accent, line width=1.8pt] coordinates {(0,0) (2,0) (2,1.6) (7,1.6) (7,6)};
          \addplot[quiet, domain=0.3:9.6, samples=40] {5.2*exp(-0.42*x)};
          \addplot[warm, domain=0.3:9.9, samples=40] {6.2*exp(-0.17*x)+0.2};
          \node[font=\scriptsize, color=accent, anchor=south] at (axis cs:1,0.05) {free};
          \node[font=\scriptsize, color=accent, anchor=south] at (axis cs:4.5,1.65) {price $e$};
          \node[font=\scriptsize, color=accent, anchor=west] at (axis cs:7.1,4.5) {cap};
          \node[font=\scriptsize, color=quiet, anchor=south] at (axis cs:8.5,0.25) {benefit, normal};
          \node[font=\scriptsize, color=warm, anchor=west] at (axis cs:2.4,5.2) {benefit, crisis};
        \end{axis}
      \end{tikzpicture}
    \end{column}
    \begin{column}{0.42\textwidth}
      \colhead{Weitzman (1974)}
      \small
      Advantage of prices over quantities:
      \[
        \Delta = \frac{\sigma^2\,(b-c)}{2\,b^2}
      \]
      \quiet{$b$: slope of benefit, $c$: slope of cost}
      \vspace{0.5em}
      \begin{itemize}
        \setlength\itemsep{0.3em}
        \item Flat cost, steep benefit: \textbf{price}
        \item Vertical cost at $L$: \textbf{quantity}
        \item Both on one curve: \textbf{tiered rate with a cap} \quiet{(Roberts \& Spence 1976)}
      \end{itemize}
    \end{column}
  \end{columns}
\end{frame}

% 9 ================================================================
\begin{frame}{The optimal digital euro has three numbers: a free tier, a Pigouvian rate above it, and a cap}{The design}
  \vspace{0.2cm}
  \begin{columns}[T]
    \begin{column}{0.29\textwidth}
      \colhead{Free tier $M_1$}
      {\fontsize{20}{22}\selectfont\color{accent}\euro 510}\par\vspace{0.2em}
      {\scriptsize $N\,\mathbb{E}[\min(t,M_1)] = X_0$}\par
      {\scriptsize\quiet{payments use only what banks absorb for free (\euro 127bn)}}
    \end{column}
    \begin{column}{0.29\textwidth}
      \colhead{Rate above the tier}
      {\fontsize{20}{22}\selectfont\color{accent}$-2.4\%$}\par\vspace{0.2em}
      {\scriptsize $-e = -(\text{MRO} - r_D)$}\par
      {\scriptsize\quiet{holders pay the bank's cost of replacing the deposit}}
    \end{column}
    \begin{column}{0.29\textwidth}
      \colhead{Cap $\bar M$}
      {\fontsize{20}{22}\selectfont\color{accent}\euro 3{,}000}\par\vspace{0.2em}
      {\scriptsize $N\,\mathbb{E}[\min(d,\bar M)] = \theta D$}\par
      {\scriptsize\quiet{crisis outflow stays within bank liquidity tolerance}}
    \end{column}
  \end{columns}
  \vspace{0.35cm}
  \footnotesize
  \begin{center}
  \begin{tabular}{@{}lccc@{}}
    & Users blocked & Outflow if rates $<0$ & Outflow in a crisis \\
    \arrayrulecolor{hair}\midrule
    Cap only \quiet{(0\%, \euro 3{,}000)} & 18.5\% & \euro 648bn & \euro 648bn \\
    Price only \quiet{(tier, $-2.4\%$)} & 0\% & \euro 158bn & \textcolor{warm}{\euro 4{,}716bn} \\
    \textbf{\hl{Hybrid}} \quiet{(tier, $-2.4\%$, cap)} & \textbf{\hl{7.3\%}} & \textbf{\hl{\euro 158bn}} & \textbf{\hl{\euro 648bn}} \\
  \end{tabular}
  \end{center}
  \source{Calibrated to ECB WP 2980 survey distributions and ECB holding-limit outflows (Oct 2025); MRO 2.65\%, household overnight rate 0.26\%.}
\end{frame}

% 10 ===============================================================
\begin{frame}{The digital euro should bear interest: zero for payments, decay above, and a cap that moves with the cycle}{Conclusion}
  \begin{columns}[c]
    \begin{column}{0.50\textwidth}
      \begin{enumerate}
        \setlength\itemsep{0.5em}
        \item Deadlines get money spent, but expiry belongs to a voucher, not to money
        \item Decay keeps every euro equal and is simply a negative rate
        \item A zero-rate digital euro is costless cash: it raises the lower bound
        \item Free tier, Pigouvian decay, and a cap tied to bank liquidity
      \end{enumerate}
    \end{column}
    \begin{column}{0.44\textwidth}
      \begin{tikzpicture}
        \begin{axis}[quietaxis, width=\textwidth, height=4.6cm, xmin=125, xmax=190, ymin=0, ymax=6000,
            xtick={130,145,157,170,185}, xticklabels={130\%,145\%,157\%,170\%,185\%},
            ytick={1000,3000,5000}, yticklabels={\euro 1k,\euro 3k,\euro 5k},
            xlabel={banks' liquidity coverage ratio}, title={Countercyclical cap}]
          \addplot[accent, mark=*, mark size=1.6pt] coordinates {(130,1278) (145,2171) (157,3000) (170,4098) (185,5560)};
          \node[font=\scriptsize, color=quiet, anchor=north west] at (axis cs:157,2950) {today};
        \end{axis}
      \end{tikzpicture}
    \end{column}
  \end{columns}
\end{frame}

% ================= BACKUP =================
\appendix
\begin{frame}[plain]
  \vfill\centering{\Large\color{quiet} Backup}\vfill
\end{frame}

\begin{frame}{Expiry versus decay: derivation}{Backup}
  \footnotesize
  \textbf{Holder's problem.} Per unit of balance, with decay $\lambda$:
  \[
    (r+\lambda)V = \max_{\kappa}\Bigl\{\kappa(1-V) - \frac{\kappa^2}{2a}\Bigr\}
    \;\Rightarrow\; \kappa^*=a(1-V),\quad (r+\lambda)V=\tfrac a2(1-V)^2 .
  \]
  \textbf{Square-root rule.} With $x=1-V$ small, $(r+\lambda)(1-x)=\tfrac a2 x^2$, so $x\approx\sqrt{2(r+\lambda)/a}$ and $\kappa^*\approx\sqrt{2a(r+\lambda)}$.

  \textbf{Expiry.} With $\tau=T-t$ left: $dV/d\tau=-rV+\tfrac a2(1-V)^2$, $V(0)=\phi$. If $\tfrac a2(1-\phi)^2>r\phi$, $V>\phi$ before $T$, so the fire sale happens only at $T$.
  Velocity rises to $a(1-\phi)$: finite, but the unspent balance is sold at once.

  \textbf{Calibration.} $\lambda$ matches the average expiry velocity over $[0,T]$.

  \vspace{0.4em}
  \begin{tabular}{@{}lcccccc@{}}
    $T$ & useful (expiry) & dumped & value range & $\lambda$ p.a. & useful (decay) & decayed \\
    \arrayrulecolor{hair}\midrule
    7 d  & 45\% & 55\% & 0.50--0.63 & 1172\% & 41\% & 15\% \\
    30 d & 84\% & 16\% & 0.50--0.80 & 489\%  & 73\% & 16\% \\
    90 d & 97\% & 3\%  & 0.50--0.91 & 161\%  & 88\% & 10\% \\
  \end{tabular}
\end{frame}

\begin{frame}{Prices versus quantities: derivation}{Backup}
  \footnotesize
  Benefit $MB(q)=b_0+\theta-bq$, cost $MC(q)=c_0+cq$, shock $\theta$ with variance $\sigma^2$; optimum $q^*(\theta)=(b_0+\theta-c_0)/(b+c)$.
  \begin{itemize}
    \setlength\itemsep{0.4em}
    \item Quantity fixed at $\mathbb E q^*$: expected loss $\dfrac{\sigma^2}{2(b+c)}$
    \item Price fixed: $q$ moves by $\theta/b$, expected loss $\dfrac{\sigma^2 c^2}{2b^2(b+c)}$
    \item Difference: $\Delta=\dfrac{\sigma^2(b-c)}{2b^2}$. Prices win iff $b>c$.
  \end{itemize}
  \vspace{0.5em}
  \textbf{Three-segment cost.} $MC(X)=0$ for $X\le X_0$, $e$ for $X_0<X\le L$, $\infty$ beyond $L$.
  The planner's optimum is decentralised by a zero rate up to $M_1$, the rate $-e$ above it, and a cap at $\bar M$:
  \[
    N\,\mathbb E[\min(t,M_1)]=X_0,\qquad i_2=-e,\qquad N\,\mathbb E[\min(d,\bar M)]=\theta D .
  \]
  Tier-2 demand shrinks by $K=\sqrt{r_D/(r_D+e)}=0.31$ (Baumol--Tobin).
\end{frame}

\begin{frame}{Calibration: data and fit}{Backup}
  \footnotesize
  \begin{tabular}{@{}lll@{}}
    Sight deposits per person & median \euro 2{,}000, mean \euro 13{,}100 & ECB WP 2980 \\
    Desired digital-euro holdings & median \euro 560, mean \euro 3{,}220 & ECB WP 2980 \\
    Crisis outflow & \euro 156bn at \euro 500; \euro 699bn at \euro 3{,}000 & ECB, Oct 2025 \\
    Digitalisation inflow by 2034 & \euro 127bn & ECB, Oct 2025 \\
    MRO; household overnight rate & 2.65\%; 0.26\% & ECB, Sep / Apr 2026 \\
    Aggregate LCR & 156.7\% & ECB, Q3 2025 \\
    \arrayrulecolor{hair}\midrule
  \end{tabular}

  \vspace{0.6em}
  \textbf{Fit.} Lognormals from the survey; population fitted to the two outflows: $N=360$m (euro area 2026: 359.6m).
  Model: \euro 156bn at \euro 500, \euro 648bn at \euro 3{,}000; mean holding under a \euro 3{,}000 cap \euro 1{,}086 (survey: \euro 1{,}100); 18\% want more than \euro 3{,}000 (survey: 20\%).
  No cap: \euro 4.7tn (Meller \& Soons 2023: almost \euro 5tn).

  \vspace{0.6em}
  \textbf{Sensitivity.} $X_0$ of \euro 100bn / 200bn / 300bn gives $M_1$ of \euro 370 / 960 / 1{,}840.
  Tolerance $\theta$ of 5\% / 7.6\% / 10\% gives $\bar M$ of \euro 1{,}690 / 3{,}000 / 4{,}520.

  \vspace{0.6em}
  \textbf{Caveat.} $e$ is the private cost of replacing deposits; with central-bank pass-through funding the social cost may be lower.
\end{frame}

\begin{frame}{References}{Backup}
  \scriptsize
  \begin{itemize}
    \setlength\itemsep{0.1em}
    \item Agarwal, R.\ \& Kimball, M.\ (2019). Enabling deep negative rates to fight recessions. IMF WP 19/84.
    \item Bindseil, U.\ (2020). Tiered CBDC and the financial system. ECB WP 2351.
    \item Boehm, J., Fize, E.\ \& Jaravel, X.\ (2024). Five facts about MPCs: evidence from a randomized experiment. CEP DP 1998.
    \item Bordo, M.\ \& Levin, A.\ (2017). Central bank digital currency and the future of monetary policy. NBER WP 23711.
    \item Brunnermeier, M.\ \& Niepelt, D.\ (2019). On the equivalence of private and public money. \emph{JME} 106.
    \item Burlon, L., Mu\~noz, M.\ \& Smets, F.\ (2024). The optimal quantity of CBDC in a bank-based economy. \emph{AEJ: Macro} 16(4).
    \item Ding, J., Jiang, L., Msall, L.\ \& Notowidigdo, M.\ (2024). Consumer-financed fiscal stimulus: evidence from digital coupons in China.
    \item European Central Bank (2025). Technical data on financial stability impact of digital euro, 9 October.
    \item European Commission (2023). Proposal for a Regulation on the establishment of the digital euro, Art.\ 16(8).
    \item Gesell, S.\ (1916). \emph{The Natural Economic Order}. Goodfriend, M.\ (2000). Overcoming the zero bound on interest rate policy. \emph{JMCB} 32(4).
    \item Kahn, C., van Oordt, M.\ \& Zhu, Y.\ (2026). Best before? Expiring central bank digital currency and loss recovery. \emph{JMCB}.
    \item Lambert, C., Larkou, C., Pancaro, C., Pellicani, A.\ \& Sintonen, M.\ (2024). Digital euro demand. ECB WP 2980.
    \item Meller, B.\ \& Soons, O.\ (2023). Know your (holding) limits. ECB Occasional Paper 326.
    \item Roberts, M.\ \& Spence, M.\ (1976). Effluent charges and licenses under uncertainty. \emph{JPubE} 5.
    \item Weitzman, M.\ (1974). Prices vs.\ quantities. \emph{Review of Economic Studies} 41(4).
    \item Xing, J., Zou, E., Yin, Z., Wang, Y.\ \& Li, Z.\ (2023). ``Quick response'' economic stimulus. \emph{AEJ: Macro} 15(4).
  \end{itemize}
\end{frame}

\end{document}
